\chapter*{Daftar Notasi}
\phantomsection
\addcontentsline{toc}{chapter}{DAFTAR NOTASI}
% Arti tiap simbol diambil dari blok "keterangan" di bawah persamaan Bab II;
% urutannya mengikuti kemunculan pertama di naskah. Memakai longtable agar
% daftar ini boleh berlanjut ke halaman berikutnya.
\begingroup
\setstretch{1.2}
\begin{longtable}{@{}p{3cm}@{}p{0.6cm}@{}p{\dimexpr\textwidth-3.6cm\relax}@{}}
  $r_{t,i}$ & : & \textit{Return} logaritmik intraday pada bar ke-$i$ hari $t$ \\
  $C_{t,i}$ & : & Harga penutupan bar 5 menit ke-$i$ pada hari $t$; untuk
    $i = 1$, $C_{t,0}$ adalah harga penutupan bar terakhir hari $t-1$ \\
  $i$ & : & Nomor urut bar 5 menit dalam satu hari, bernilai 1 sampai 288 \\
  $RV_t$ & : & \textit{Realized variance} harian pada hari $t$, yaitu
    komponen harian \\
  $RV_t^{(w)}$ & : & Rata-rata RV tujuh hari terakhir hingga hari $t$, yaitu
    komponen mingguan \\
  $RV_t^{(m)}$ & : & Rata-rata RV tiga puluh hari terakhir hingga hari $t$,
    yaitu komponen bulanan \\
  $j$ & : & Indeks pergeseran hari ke belakang dari hari $t$ \\
  $RV_{t+1}$ & : & RV hari berikutnya, yaitu target peramalan \\
  $\beta_0$ & : & Konstanta regresi HAR-RV \\
  $\beta_d$ & : & Koefisien komponen harian \\
  $\beta_w$ & : & Koefisien komponen mingguan \\
  $\beta_m$ & : & Koefisien komponen bulanan \\
  $\varepsilon_{t+1}$ & : & Galat regresi HAR-RV yang diasumsikan
    berrata-rata nol \\
  $F_b(\mathbf{x})$ & : & Model GBDT gabungan setelah $b$ iterasi \\
  $h_b(\mathbf{x})$ & : & Pohon ke-$b$, dilatih pada gradien fungsi
    \textit{loss} \\
  $\eta$ & : & Laju pembelajaran \\
  $RS_t^-$ & : & \textit{Realized semivariance} negatif pada hari $t$ \\
  $\mathbf{1}(\cdot)$ & : & Fungsi indikator, bernilai 1 bila syarat di
    dalamnya terpenuhi dan 0 bila tidak \\
  $\sigma^2_{P,t}$ & : & Estimator varians \textit{range} Parkinson pada
    hari $t$ \\
  $P^{H}_t$ & : & Harga tertinggi pada hari $t$ \\
  $P^{L}_t$ & : & Harga terendah pada hari $t$ \\
  $\text{Amihud}_t$ & : & Ukuran \textit{illiquidity} pada hari $t$ \\
  $QV_{t,i}$ & : & Volume dalam dolar (USDT) pada bar ke-$i$ \\
  $N_t$ & : & Jumlah bar bervolume positif pada hari $t$ \\
  $10^6$ & : & Faktor penskalaan ukuran Amihud agar besarannya mudah dibaca \\
  $F_t$ & : & Ramalan RV untuk hari $t$ \\
  $L^{\text{QLIKE}}_t$ & : & Nilai \textit{loss} QLIKE pada hari $t$ \\
  $L^{\text{MSE}}_t$ & : & Nilai \textit{loss} MSE pada hari $t$ \\
  $\Delta L_t$ & : & Selisih \textit{loss} harian antara model
    \textit{machine learning} dan HAR-RV \\
  $L_t(\text{ML})$ & : & \textit{Loss} model \textit{machine learning} pada
    hari $t$ \\
  $L_t(\text{HAR})$ & : & \textit{Loss} model HAR-RV pada hari $t$ \\
  $DM$ & : & Statistik uji Diebold--Mariano \\
  $\overline{\Delta L}$ & : & Rata-rata selisih \textit{loss} selama periode
    uji \\
  $\widehat{V}_{\text{HAC}}$ & : & Estimator varians jangka panjang yang
    konsisten terhadap heteroskedastisitas dan autokorelasi \\
  $n$ & : & Jumlah observasi pada periode uji \\
  $\text{UDmax}$ & : & Statistik uji perubahan struktural maksimum \\
  $\sup F_T(m)$ & : & Statistik $F$ terbesar atas semua kemungkinan
    penempatan $m$ titik patah \\
  $m$ & : & Jumlah titik patah pada uji Bai--Perron \\
  $M$ & : & Jumlah maksimum titik patah yang diuji \\
  $\varepsilon$ & : & Proporsi \textit{trimming}, yaitu panjang minimum tiap
    segmen pada uji Bai--Perron \\
  $\mu$ & : & Fraksi ukuran jendela bergulir pada \textit{Fluctuation test},
    dengan $n_F = \mu n$ \\
  $n_F$ & : & Ukuran jendela bergulir pada \textit{Fluctuation test} \\
  $D^{\text{KS}}$ & : & Statistik Kolmogorov--Smirnov dua sampel \\
  $\hat{F}_{\text{latih}}(x)$ & : & Fungsi distribusi kumulatif empiris data
    latih \\
  $\hat{F}_{\text{uji}}(x)$ & : & Fungsi distribusi kumulatif empiris bulan
    uji \\
  $\sup_x$ & : & Nilai terbesar atas seluruh $x$ \\
  $D_t$ & : & Selisih $\Delta L_t$ antara strategi pembaruan S2-365 dan S1 \\
  $W$ & : & Panjang jendela data latih \textit{rolling}, dalam hari \\
  $C_t$ & : & Harga penutupan bar terakhir hari $t$ \\
  $t^*$ & : & Hari contoh pada perhitungan manual (30 Agustus 2018) \\
  $\hat{y}_{t+1}$ & : & Ramalan $\ln RV_{t+1}$ sebelum dikembalikan ke skala
    level \\
  $\hat{\sigma}^2$ & : & Varians residual untuk koreksi bias transformasi
    logaritma \\
  $\underline{F}$ & : & Batas bawah ramalan, persentil ke-1 target $RV$ pada
    \textit{window} latih \\
  $F_{t+1}$ & : & Ramalan RV hari $t+1$ setelah koreksi bias dan batas bawah \\
  $\bar{p}_t$ & : & Rata-rata logaritma harga tertinggi dan terendah hari
    $t$ pada estimator Abdi--Ranaldo \\
  $c_t$ & : & Logaritma harga penutupan hari $t$ \\
  $s_t$ & : & \textit{Spread} Abdi--Ranaldo harian \\
  $\tau$ & : & Indeks bulan uji pada regresi karakteristik pasar \\
  $\overline{\Delta L}_\tau$ & : & Rata-rata $\Delta L_t$ pada bulan $\tau$ \\
  $z_{k,\tau}$ & : & Variabel karakteristik pasar ke-$k$ pada bulan $\tau$ \\
  $k$ & : & Indeks variabel karakteristik pasar, bernilai 1 sampai 4 \\
  $\gamma_k$ & : & Koefisien variabel karakteristik pasar ke-$k$ \\
  $\alpha$ & : & Konstanta regresi karakteristik pasar dan regresi era \\
  $\delta$ & : & Koefisien variabel \textit{dummy} era \\
  $\text{Era}_\tau$ & : & Variabel \textit{dummy} era \textit{institutional}
    pada bulan $\tau$ \\
  $\text{Era}_t$ & : & Variabel \textit{dummy} era \textit{institutional}
    pada hari $t$ \\
  $u_\tau$ & : & Galat regresi karakteristik pasar bulan $\tau$ \\
  $e_t$ & : & Galat regresi era pada hari $t$ \\
  $\ell$ & : & Jumlah titik patah pada uji perubahan berurutan
    $\sup F(\ell+1 \mid \ell)$ \\
  $R$ & : & Jumlah replikasi simulasi Monte Carlo \\
  $R^{+}$ & : & Jumlah statistik simulasi yang tidak lebih kecil daripada
    statistik amatan \\
\end{longtable}
\endgroup
\clearpage
